\documentclass[11pt,a4paper]{article}
% preamble.tex — estilo común de todos los apuntes Froth.
% Cada apunte hace % preamble.tex — estilo común de todos los apuntes Froth.
% Cada apunte hace % preamble.tex — estilo común de todos los apuntes Froth.
% Cada apunte hace \input{preamble} justo después de \documentclass.
\usepackage[margin=2.4cm]{geometry}
\usepackage{xcolor}
\definecolor{litblue}{RGB}{31,79,121}
\definecolor{litgreen}{RGB}{27,94,32}
\definecolor{litorange}{RGB}{230,81,0}
\usepackage{parskip}
\usepackage{titlesec}
\titleformat{\section}{\Large\bfseries\color{litblue}}{\thesection}{0.6em}{}
\titleformat{\subsection}{\large\bfseries\color{litblue}}{\thesubsection}{0.6em}{}
\usepackage{tcolorbox}
\usepackage{fancyhdr}
\pagestyle{fancy}
\fancyhf{}
\fancyhead[L]{\small\color{litblue}Froth — Apuntes de ML}
\fancyhead[R]{\small\thepage}
\renewcommand{\headrulewidth}{0.4pt}
\usepackage[colorlinks=true,linkcolor=litblue,urlcolor=litblue]{hyperref}

% Cabecera de cada apunte: \cabecera{numero}{titulo}
\newcommand{\cabecera}[2]{%
  \begin{tcolorbox}[colback=litblue,coltext=white,colframe=litblue,arc=2mm]
    {\Large\bfseries Apunte #1 — #2}\\[3pt]
    {\small Proyecto Froth · aprender Machine Learning construyendo}
  \end{tcolorbox}\vspace{0.8em}}

% Cajas temáticas
\newtcolorbox{concepto}[1]{colback=blue!4,colframe=litblue,title={Concepto: #1},arc=1.5mm}
\newtcolorbox{practica}{colback=green!5,colframe=litgreen,title={Pruébalo tú (teclado en mano)},arc=1.5mm}
\newtcolorbox{ojo}{colback=orange!8,colframe=litorange,title={Ojo / error típico},arc=1.5mm}
\newtcolorbox{resumen}{colback=gray!8,colframe=gray!60,title={En una frase},arc=1.5mm}

\newcommand{\codigo}[1]{\texttt{#1}}
 justo después de \documentclass.
\usepackage[margin=2.4cm]{geometry}
\usepackage{xcolor}
\definecolor{litblue}{RGB}{31,79,121}
\definecolor{litgreen}{RGB}{27,94,32}
\definecolor{litorange}{RGB}{230,81,0}
\usepackage{parskip}
\usepackage{titlesec}
\titleformat{\section}{\Large\bfseries\color{litblue}}{\thesection}{0.6em}{}
\titleformat{\subsection}{\large\bfseries\color{litblue}}{\thesubsection}{0.6em}{}
\usepackage{tcolorbox}
\usepackage{fancyhdr}
\pagestyle{fancy}
\fancyhf{}
\fancyhead[L]{\small\color{litblue}Froth — Apuntes de ML}
\fancyhead[R]{\small\thepage}
\renewcommand{\headrulewidth}{0.4pt}
\usepackage[colorlinks=true,linkcolor=litblue,urlcolor=litblue]{hyperref}

% Cabecera de cada apunte: \cabecera{numero}{titulo}
\newcommand{\cabecera}[2]{%
  \begin{tcolorbox}[colback=litblue,coltext=white,colframe=litblue,arc=2mm]
    {\Large\bfseries Apunte #1 — #2}\\[3pt]
    {\small Proyecto Froth · aprender Machine Learning construyendo}
  \end{tcolorbox}\vspace{0.8em}}

% Cajas temáticas
\newtcolorbox{concepto}[1]{colback=blue!4,colframe=litblue,title={Concepto: #1},arc=1.5mm}
\newtcolorbox{practica}{colback=green!5,colframe=litgreen,title={Pruébalo tú (teclado en mano)},arc=1.5mm}
\newtcolorbox{ojo}{colback=orange!8,colframe=litorange,title={Ojo / error típico},arc=1.5mm}
\newtcolorbox{resumen}{colback=gray!8,colframe=gray!60,title={En una frase},arc=1.5mm}

\newcommand{\codigo}[1]{\texttt{#1}}
 justo después de \documentclass.
\usepackage[margin=2.4cm]{geometry}
\usepackage{xcolor}
\definecolor{litblue}{RGB}{31,79,121}
\definecolor{litgreen}{RGB}{27,94,32}
\definecolor{litorange}{RGB}{230,81,0}
\usepackage{parskip}
\usepackage{titlesec}
\titleformat{\section}{\Large\bfseries\color{litblue}}{\thesection}{0.6em}{}
\titleformat{\subsection}{\large\bfseries\color{litblue}}{\thesubsection}{0.6em}{}
\usepackage{tcolorbox}
\usepackage{fancyhdr}
\pagestyle{fancy}
\fancyhf{}
\fancyhead[L]{\small\color{litblue}Froth — Apuntes de ML}
\fancyhead[R]{\small\thepage}
\renewcommand{\headrulewidth}{0.4pt}
\usepackage[colorlinks=true,linkcolor=litblue,urlcolor=litblue]{hyperref}

% Cabecera de cada apunte: \cabecera{numero}{titulo}
\newcommand{\cabecera}[2]{%
  \begin{tcolorbox}[colback=litblue,coltext=white,colframe=litblue,arc=2mm]
    {\Large\bfseries Apunte #1 — #2}\\[3pt]
    {\small Proyecto Froth · aprender Machine Learning construyendo}
  \end{tcolorbox}\vspace{0.8em}}

% Cajas temáticas
\newtcolorbox{concepto}[1]{colback=blue!4,colframe=litblue,title={Concepto: #1},arc=1.5mm}
\newtcolorbox{practica}{colback=green!5,colframe=litgreen,title={Pruébalo tú (teclado en mano)},arc=1.5mm}
\newtcolorbox{ojo}{colback=orange!8,colframe=litorange,title={Ojo / error típico},arc=1.5mm}
\newtcolorbox{resumen}{colback=gray!8,colframe=gray!60,title={En una frase},arc=1.5mm}

\newcommand{\codigo}[1]{\texttt{#1}}

\usepackage{tikz}
\begin{document}
\cabecera{19}{Grafos: nodos, aristas, hubs y puentes}

\section{Qué es un grafo}

\begin{concepto}{grafo = cosas + relaciones}
Piensa en tu red social: personas $=$ \textbf{nodos}, amistades $=$ \textbf{aristas} (las
líneas). Esa estructura —cosas conectadas por relaciones— es un \textbf{grafo}, y está detrás
de Google Maps (cruces + calles), LinkedIn (gente + contactos) y nuestro mapa mental
(papers + ``se relacionan con''). La librería estándar en Python: \codigo{networkx}.
\end{concepto}

Cuatro ideas que usaremos:
\begin{enumerate}
  \item \textbf{Peso}: no todas las relaciones son igual de fuertes. ``Se parecen 0.94'' pesa
        más que ``se parecen 0.60'' — la línea se dibuja más gruesa.
  \item \textbf{Dirección}: la similitud es mutua (arista sin flecha); las citas no — A cita
        a B, no al revés (arista \emph{con} flecha).
  \item \textbf{Hub}: nodo con muchísimas conexiones.
  \item \textbf{Puente}: nodo o arista que conecta dos comunidades que casi no se tocan.
\end{enumerate}

\begin{center}
\begin{tikzpicture}[scale=1.0, every node/.style={circle, inner sep=2pt}]
  % community A (flotation, green-ish left)
  \foreach \p/\n in {(0,1)/a1, (0.9,1.7)/a2, (1.0,0.4)/a3, (1.8,1.1)/a4}
    \fill[litgreen!70] \p circle (5pt) coordinate (\n);
  \draw[litgreen!70] (a1)--(a2) (a1)--(a3) (a2)--(a4) (a3)--(a4) (a1)--(a4) (a2)--(a3);
  \node[litgreen] at (0.9,2.3) {\small comunidad A (flotación)};
  % community B (geology, blue right)
  \foreach \p/\n in {(5.6,1.0)/b1, (6.5,1.7)/b2, (6.6,0.3)/b3, (7.4,1.0)/b4}
    \fill[litblue!70] \p circle (5pt) coordinate (\n);
  \draw[litblue!70] (b1)--(b2) (b1)--(b3) (b2)--(b4) (b3)--(b4) (b1)--(b4) (b2)--(b3);
  \node[litblue] at (6.5,2.3) {\small comunidad B (geología)};
  % bridge node
  \fill[litorange] (3.7,0.9) circle (6pt) coordinate (br);
  \draw[litorange, thick] (a4)--(br)--(b1);
  \node[litorange] at (3.7,0.2) {\small \textbf{puente}};
  % hub marker: a4 has most connections
  \node[litgreen] at (2.45,1.55) {\small hub};
\end{tikzpicture}
\end{center}

\section{Para qué sirve cada uno en TU literature review}

\begin{concepto}{hubs = cimiento; puentes = originalidad}
\begin{itemize}
  \item \textbf{El hub} es la lectura obligatoria: el clásico que todos citan. Va en la
        introducción de tu review — demuestra que conoces el terreno. Pero \emph{todo el
        mundo conoce los hubs}: ahí no hay novedad posible.
  \item \textbf{El puente} conecta dos temas que casi no se hablan. Si solo 2 papers unen
        ``flotación'' con ``geología del yacimiento'', esa frontera está poco explorada
        $\rightarrow$ ahí cabe una contribución nueva. Los \emph{gaps} de la Fase 5 son,
        literalmente, puentes que faltan.
\end{itemize}
Respuesta corta: hubs para fundamentar, \textbf{puentes para aportar}.
\end{concepto}

\begin{ojo}
\textbf{Tu tesis ES un puente.} Mira el título: \emph{``Flotation of ... Li bearing mica
minerals \textbf{and its impact on} the by-products (Ta, Nb, Be,...) recovery...''}. Conecta
exactamente el cluster verde (flotación) con el azul (geología/by-products del granito).
No es un tema \emph{dentro} de una comunidad: es la \emph{frontera} entre dos. Por eso es
una tesis y no una tarea. Cuando la red esté construida, búscate en ella.
\end{ojo}

\section{Qué grafo construiremos (lo que viene)}

\begin{enumerate}
  \item \textbf{Red de similitud} (4.2): cada paper conectado con sus k vecinos más parecidos
        por coseno \textbf{en 768D} (el espacio real, no el mapa 2D — Apunte 12). Con umbral,
        para que no salga una ``bola de pelo''.
  \item \textbf{Red de citas} (4.3): flechas de quién cita a quién. Requiere re-bajar los
        papers pidiendo el campo \codigo{referenced\_works} que en la Fase 1 no pedimos.
  \item \textbf{Mapa mental de temas} (4.4): un nodo por subtema, aristas $=$ qué tan
        relacionados están. La vista presentable.
\end{enumerate}

\begin{resumen}
Un grafo son nodos + aristas (con peso y a veces dirección). Los hubs son los clásicos que
fundamentan; los puentes, las fronteras donde vive la originalidad — y tu propia tesis es un
puente entre flotación y geología. Construiremos redes de similitud, de citas y de temas.
\end{resumen}

\end{document}
